\uuid{Vbh8}
\titre{ Initialisation des poids : analyse de trois stratégies }
\niveau{L3}
\module{Analyse de données}
\chapitre{Réseaux de neurones}
\sousChapitre{Initialisation, symétrie, explosion du gradient}
\theme{Initialisation à zéro, initialisation aléatoire, brisure de symétrie}
\auteur{Maxime Nguyen}
\datecreate{2026-03-19}
\organisation{AMSCC}
\difficulte{2}

\contenu{
\texte{
  On considère trois initialisations pour une matrice de poids $W \in \mathbb{R}^{4 \times 3}$ :
  \begin{itemize}
    \item[A.] \texttt{W = np.zeros((4, 3))}
    \item[B.] \texttt{W = np.random.randn(4, 3) * 0.01}
    \item[C.] \texttt{W = np.random.randn(4, 3) * 100}
  \end{itemize}
}
\begin{enumerate}

\item
  \question{Laquelle de ces initialisations est problématique, et pourquoi ?}
  \reponse{
    \textbf{A} est problématique : initialiser tous les poids à zéro brise la rétropropagation par \emph{symétrie}.
    Tous les neurones d'une même couche calculent la même sortie et reçoivent les mêmes gradients,
    donc ils évoluent identiquement : la couche cachée est équivalente à un seul neurone, quelle que soit sa largeur.

    \textbf{C} est également problématique : des poids initiaux très grands ($\times 100$) provoquent
    une saturation des activations ou une explosion des gradients dès le début de l'entraînement.

    \textbf{B} est la bonne pratique : petite valeur aléatoire, brise la symétrie sans saturer les activations.
  }

\end{enumerate}
}
